\documentclass[11pt,letterpaper]{article}
\usepackage{bes-whitepaper}

\begin{document}

\wptitle{Losing iron sites or losing the nickel host? Why Ni--Fe oxygen-evolution\\catalysts deactivate}
  {Catalysis Science}{Viviane Schwartz, Chris Bradley}

\hd{Research question.} Ni--Fe (oxy)hydroxides are among the most active earth-abundant catalysts for
the oxygen evolution reaction (OER) in alkaline media [1,2], the anode chemistry of alkaline and
anion-exchange-membrane electrolyzers, which need no iridium. Their activity fades under load. Is
that because Fe active sites leave faster than they return, or because the Ni host itself
transforms?

\hd{Knowledge gap.} Fe in these catalysts is not fixed: it dissolves and redeposits during OER, and
activity holds only while the two rates balance [3]. Chung et al.\ showed this across several hosts
and proposed a descriptor for Fe retention [3]. Whether that exchange, or restructuring of the host
such as Ni/Fe segregation, sets durability when both are measured as competing rate laws on the same
films has not been tested, because each test consumes a film and needs trace electrolyte analysis.

\hd{Hypotheses.} \emph{H1, site exchange:} retained activity is set by the steady-state Fe site
coverage, fixed by a composition-dependent Fe dissolution rate and an electrolyte-Fe-dependent
redeposition rate. Predictions: adding Fe to the electrolyte restores activity, and Fe lost to the
electrolyte predicts the activity decay of every composition with the same two rate laws.
\emph{H2, host restructuring:} activity loss follows irreversible change in the Ni host (fewer
redox-active Ni sites, rising charge-transfer resistance, loss of film). Predictions: added Fe does
not restore activity, and decay tracks Ni lost from the film (digested after the test) and a change
of phase, not the Fe flux.

\hd{Approach.} An Opentrons OT-2 fitted with the open AMPERE-2 electrodeposition hardware [4]
plates matched Ni--Fe--M films (Fe 0--25\%; M = Co or Mn, 0--20\%) on identical substrates by one
fixed protocol, then conditions them anodically: the films are a controlled variable, not the product. Each film gets CV, a Tafel
scan and EIS, then a short stress test (constant current or potential cycling) in Fe-free KOH [2]
with Fe added at set levels, while the robot draws electrolyte aliquots for ICP-MS (Fe, Ni, M); an $^{57}$Fe spike separates
redeposition from dissolution. Both
mechanisms are written as kinetic models with shared parameters. Bayesian experimental design picks
the next composition, electrolyte Fe level and current density where the models' predicted activity
retention differs most, so AI serves mechanism identification.

\wpflow
  {\textbf{Design}\\ next film, Fe level and current where H1 and H2 disagree}
  {\textbf{OT-2 + AMPERE-2}\\ plate a matched Ni--Fe--M film; set electrolyte Fe}
  {\textbf{Electrochemistry}\\ CV, Tafel, EIS; stress test; ICP-MS, $^{57}$Fe}
  {\textbf{Kinetic models}\\ fit site-exchange (H1) and host-loss (H2)}
  {\textbf{Outcome}\\ which mechanism controls where; Fe exchange rates}
  {update the probability of H1 and H2, then choose the next film}

\hd{Expected outcomes.} (1) A map of where in composition and electrolyte space each mechanism
controls durability; (2) Fe dissolution and redeposition rate constants as functions of host
composition; (3) whether a third metal works by slowing Fe loss or by stabilizing the host. Together
these explain how to keep earth-abundant active sites in place over an electrolyzer's life.

\hd{Fit and resources.} The project addresses catalyst \emph{dynamics during catalytic cycles and
deactivation}, \emph{mechanisms mediated by earth-abundant metals}, and
\emph{AI/ML for mechanism identification}. It does not optimize catalyst synthesis or electrolyzer
design. The PI works on self-driving labs and Bayesian experimental design; the lab has an OT-2, is
adding a potentiostat, and has campus ICP access.

\begin{wprefs}
[1] M.\,W. Louie, A.\,T. Bell, \emph{J.\ Am.\ Chem.\ Soc.} 135, 12329 (2013), doi:10.1021/ja405351s.
[2] L. Trotochaud \emph{et al.}, \emph{J.\ Am.\ Chem.\ Soc.} 136, 6744 (2014), doi:10.1021/ja502379c.
[3] D.\,Y. Chung \emph{et al.}, \emph{Nat.\ Energy} 5, 222 (2020), doi:10.1038/s41560-020-0576-y.
[4] Fisker-B{\o}dker \emph{et al.}, \emph{Digital Discovery} 4, 2491 (2025), doi:10.1039/D5DD00180C.
\end{wprefs}

\end{document}
